
\section{conclusion}

\begin{frame}
\frametitle{Multiphoton Ionization Reduction of Atoms in Two-Color
Femtosecond Laser Fields \small (Phys. Rev. Lett. \textbf{130}, 113201 (2023))}
\framesubtitle{Conclusion}


\small
\begin{textblock*}{400pt}(30pt,60pt)
\vspace{30pt}
In this paper, the mechanism of ionization reduction 
in strong field is explorerd
\begin{itemize}
\footnotesize
\item For atoms exposed to SH-FW laser fields, 
the weaker FW shunts the wavepacket to higher Rydberg states, 
though more loosely bound, are harder to ionize
\item By tuning the intensity of the FW, the populations 
of the Rydberg and continuum states can be controlled
%\begin{itemize}
%\footnotesize
%\item[$\bullet$] Rare quantitative information can be probed about 
%various high intensity effects
%\item[$\bullet$] Importance of long range electronic potential 
%and nonadiabatic effects are better rendered 
%\end{itemize}

\vspace{20pt}

\item So this is yet another study shedding light on ionization rate control.
\end{itemize}

\end{textblock*}


%\begin{textblock*}{100pt}(260pt,70pt)
%\includegraphics[scale=.30]{images/medium1.png}
%\end{textblock*}

\end{frame}
